% Do not change document class, margins, fonts, etc.
\documentclass[a4paper,oneside,bibliography=totoc]{scrbook}

% some useful packages (add more as needed)
\usepackage[utf8]{inputenc}
\usepackage{graphicx}
\usepackage{latexsym}
\usepackage{amsmath}
\usepackage{amssymb}
\usepackage{tabularx}
\usepackage{booktabs}
\usepackage{algorithm}
\counterwithin{algorithm}{chapter}
\usepackage{algorithmic}
\usepackage{csquotes}
\renewcommand{\algorithmiccomment}[1]{\hfill\textit{// #1}}
\usepackage[usenames,dvipsnames]{xcolor}
\usepackage{soul}
\sethlcolor{yellow}
\usepackage[colorlinks,citecolor=Green]{hyperref}
\usepackage{enumitem}
\setlist[itemize]{noitemsep, topsep=2pt}
\setlist[enumerate]{noitemsep, topsep=2pt}

% Remove running headers, use plain page numbers only
\usepackage{scrlayer-scrpage}
\pagestyle{plain}

% Prevent chapters from forcing a new page, while keeping KOMA heading style
\makeatletter
\renewcommand\chapter{\if@openright\else\fi
  \thispagestyle{plain}%
  \global\@topnum\z@
  \@afterindentfalse
  \secdef\@chapter\@schapter}
\makeatother

% IEEE citation style — IEEEtrantools provides the bbl commands outside the IEEEtran class
\usepackage{IEEEtrantools}
\bibliographystyle{IEEEtran}
\usepackage{microtype}

% example environments
\newtheorem{definition}{Definition}
\newtheorem{proposition}{Proposition}

\begin{document}

\frontmatter
\subject{Project Outline: IE500 Data Mining}
\title{Ethnic Bias and Generalization in Facial Attractiveness Prediction}
\author{Team Number: 3 \\
        Maxim Froehlich, David Siregar, Timon Kuhl, \\
        Lars Bullmahn, Jagadeesh Gunti, Simon Heinz}
\date{April 12, 2026}
\publishers{{\small Submitted to}\\
  Prof.\ Dr.\ Christian Bizer\\
  University of Mannheim\\}
\maketitle

\mainmatter

\chapter{Problem Description}
\label{ch:problem}

Automated systems for rating or ranking human faces are increasingly used in applications ranging from online platforms to hiring tools. Unlike many computer vision tasks, facial attractiveness prediction is inherently subjective and culturally shaped, which makes it particularly prone to encoding demographic bias. Models trained on data from a narrow set of ethnic groups may learn cultural beauty standards rather than broadly applicable visual patterns, raising both methodological and ethical concerns. We study this problem as a supervised regression task, where the goal is to predict the mean attractiveness rating assigned by human annotators to a given face image.

Our core research questions are: (1) Does a model trained on SCUT-FBP5500 \cite{liang2018scut} (predominantly Asian/Caucasian) perform significantly worse on MEBeauty \cite{lebedeva2021mebeauty}, which includes Asian, Black, Caucasian, Hispanic, Indian, and Middle Eastern subjects? (2) Does training on the more diverse MEBeauty dataset reduce racial disparities in prediction accuracy? (3) \hl{Which geometric facial measurements (e.g., eye spacing, facial symmetry) drive attractiveness predictions across ethnic groups, and how well do interpretable landmark-based features perform relative to end-to-end deep models?}

\chapter{Data Acquisition}
\label{ch:data}

We use three publicly available datasets. \textbf{SCUT-FBP5500} \cite{liang2018scut} contains 5,500 face images (Asian and Caucasian subjects) with attractiveness ratings from 60 annotators per image on a 1--5 scale, available on Kaggle. \textbf{MEBeauty} \cite{lebedeva2021mebeauty} contains 2,550 images across six ethnic groups (Asian, Black, Caucasian, Hispanic, Indian, Middle Eastern) with ratings on a 1--10 scale, available on GitHub (fbplab/MEBeauty-database). \textbf{LiveBeauty} contains approximately 10,000 Asian face images with around 200,000 attractiveness labels, providing a large-scale Asian-focused complement to the other two corpora. We will verify completeness (expected missing rate ${<}1\%$) and check for corrupt files.

For cross-dataset evaluation, we treat the two datasets as separate corpora: models trained on one are tested on the other to quantify generalization failure. To make this comparison meaningful, we rescale both datasets' ratings to $[0,1]$ before modeling. Within MEBeauty, group-stratified splits enable per-ethnicity performance breakdowns.

\chapter{Methodology}
\label{ch:method}

\section{Preprocessing}

\begin{itemize}
    \item \textbf{Face Detection \& Alignment:} Standard face detection crops and aligns all images to a consistent resolution across both datasets.
    \item \hl{\textbf{Feature Extraction:} Interpretable tabular features computed from MediaPipe facial landmarks serve as the primary input representation, including geometric measurements such as interpupillary ratio, eye spacing, and mouth-to-chin ratio. Additional features may be engineered based on feature importance analysis.}
    \item \textbf{Normalization:} Rating scores may be rescaled or standardized depending on the modeling approach (e.g., z-normalization to zero mean and unit variance); tabular geometric features standardized using z-normalization.
    \item \textbf{Outlier Handling:} Images with high annotator variance (std~$> 1.0$) flagged; we experiment with their removal.
\end{itemize}

\section{Proposed Algorithms}

\noindent\textbf{Baseline:} Global mean predictor; Linear Regression on \hl{tabular geometric features}; Ridge Regression with regularization strength tuned via Optuna.

\noindent\textbf{Symbolic:} Decision Tree Regressor with tree depth and minimum samples per leaf tuned via Optuna --- enables inspection of which geometric features drive predictions and identification of interpretable decision rules.

\noindent\textbf{Ensemble:} Stacking ensemble of XGBoost, Random Forest, Gradient Boosting, and Ridge base models with a Ridge meta-learner (5-fold CV), with feature importance analysis across ethnic groups.

\noindent\textbf{Boosting:} XGBoost Regressor with tree depth, learning rate, and number of estimators tuned via Optuna to capture nonlinear interactions between geometric features on \hl{tabular data}.

\noindent\textbf{Neural Baseline:} A small multilayer perceptron trained on tabular features with Adam optimizer, learning rate and batch size tuned via Optuna, up to 30 epochs with early stopping on validation MSE.

Within each dataset, hyperparameters are tuned via Optuna with stratified 5-fold cross-validation and final evaluation uses an 80/20 stratified train/test split. For cross-dataset evaluation, models are trained on one dataset and tested directly on the other. Feature importance analysis will guide potential expansion of the geometric feature set.

\chapter{Evaluation and Expected Results}
\label{ch:evaluation}

\section{Measuring Success}

\noindent\textbf{Regression metrics:} MSE/RMSE (primary), Pearson Correlation Coefficient (PCC), and MAE.

\noindent\textbf{Generalization:} Train on SCUT, test on MEBeauty (and vice versa). The performance drop quantifies cross-ethnic generalization failure.

\noindent\textbf{Fairness / Bias:} Fairness is operationalized as consistency of predictive performance across demographic groups. We compute per-ethnicity RMSE within MEBeauty for each model and summarize disparity using: (1) the absolute gap between worst-group and best-group RMSE, (2) the ratio of worst-to-best RMSE (values closer to 1 indicate fairer performance), and (3) the standard deviation of per-group RMSEs. These metrics quantify whether the model systematically underperforms for specific ethnic groups, which would indicate learned bias toward certain demographic features or beauty standards.

\noindent\textbf{State-of-the-art comparison:} Results benchmarked against published baselines for SCUT-FBP5500 (prior CNN approaches achieve PCC~$\approx 0.90$) and related literature using end-to-end CNN or transformer architectures. These serve as external comparison points rather than models we necessarily reimplement.

\section{Expected Outcomes}

We expect tabular geometric features to provide interpretable and computationally efficient predictions, while showing notable cross-dataset performance drops when training on SCUT and testing on MEBeauty. We anticipate higher errors for underrepresented groups (Black, Hispanic, Middle Eastern) in SCUT-trained models and partial mitigation when training on MEBeauty. \hl{Feature importance analysis will reveal which geometric measurements are most predictive and whether patterns differ across ethnic groups, and comparison with published CNN/transformer results will position the performance-interpretability tradeoff of our approach.}

\bibliography{refs}

\end{document}
